\documentclass[11pt]{article}
\usepackage[margin=1in]{geometry}
\usepackage{xcolor}
\usepackage{fancyhdr}
\usepackage{amssymb}
\usepackage{parskip}

\definecolor{cmdbg}{gray}{0.93}
\newcommand{\cmd}[1]{\colorbox{cmdbg}{\texttt{#1}}}
\newenvironment{cmdblock}{\begin{quote}\small\begin{ttfamily}}{\end{ttfamily}\end{quote}}
\newcommand{\checkbox}{$\square$~}

\pagestyle{fancy}
\fancyhf{}
\fancyhead[L]{\small Shadow PC Trial Guide}
\fancyhead[R]{\small MSAI-LLM-PERFORMANCE}
\fancyfoot[C]{\small Page \thepage}

\title{Shadow PC Trial --- Complete Setup \& Run Guide\\
\large From a blank Shadow.tech Windows PC to committed benchmark data\\
\normalsize Repository: \texttt{github.com/ericmedlock/MSAI-LLM-PERFORMANCE} \quad Guide version: 2026-07-10}
\author{ITCS 6881 --- LLM Execution Architecture Benchmark}
\date{}

\begin{document}
\maketitle
\thispagestyle{fancy}

\section*{What this guide does}

By the end you will have: (1) a Shadow PC with all tooling installed, (2) a
completed \emph{frontier-v2 trial} --- 36 frozen benchmark tasks $\times$ 3
architectures $\times$ 1 trial $=$ \textbf{108 benchmark runs} against the pinned
DeepSeek-R1 14B model on the RTX A4500 --- and (3) the results committed and pushed
back to GitHub. This is the harness's \textbf{first-ever NVIDIA/CUDA environment},
so the trial's real purpose is validating GPU telemetry before the university HPC
sweep. Follow the steps in order; each has a verification so you always know
whether it worked.

\textbf{Budget:} $\sim$15 minutes of installs, $\sim$10 GB of downloads,
then 2--3 hours of unattended running. \textbf{Disk needed:} $\sim$15 GB free.

\subsection*{Conventions used in this guide}

\begin{itemize}
  \item Windows has two different command windows, and it matters which one you use:
    \begin{itemize}
      \item \textbf{PowerShell} --- the blue Windows shell. Open it: press the
        \textbf{Start} key, type \texttt{powershell}, press Enter.
      \item \textbf{Git Bash} --- installed in Part 2; a Linux-style shell. Open it:
        Start key, type \texttt{git bash}, Enter.
    \end{itemize}
    Every command block below is labeled \textbf{[PowerShell]} or \textbf{[Git Bash]}.
  \item Type commands exactly as printed, then press Enter. Commands are case-sensitive
    in Git Bash.
  \item \checkbox marks a verification checkpoint --- confirm it before moving on.
  \item If Windows pops a dialog asking ``Do you want to allow this app to make
    changes?'' (a UAC prompt), click \textbf{Yes}. This happens during installs; it is expected.
\end{itemize}

\section*{Part 0 --- Before you start (Shadow-specific)}

\begin{enumerate}
  \item Connect to your Shadow PC with the Shadow client and stay connected.
    \textbf{Important:} Shadow \emph{shuts the virtual machine down a few minutes after
    you disconnect}. If that happens mid-run, no data is lost (every finished run is
    saved instantly), but the benchmark stops until you reconnect and resume. Plan to
    keep the session open for the full 2--3 hours, or check in periodically.
  \item Check free disk space: open \textbf{File Explorer} $\rightarrow$ This PC
    $\rightarrow$ the C: drive should show at least 15 GB free.
  \item \checkbox You are on the Shadow desktop with a working keyboard/mouse and
    15+ GB free.
\end{enumerate}

\section*{Part 1 --- Verify the GPU (30 seconds)}

Everything depends on the NVIDIA driver being present. Shadow pre-installs it;
verify rather than assume.

\textbf{[PowerShell]}
\begin{cmdblock}
nvidia-smi
\end{cmdblock}

\checkbox You see a table whose header line includes \texttt{NVIDIA RTX A4500}
and a memory column reading about \texttt{24564MiB} total.

\textbf{If instead you see} ``\texttt{nvidia-smi is not recognized}'': the driver is
missing or PATH is stale. Reboot the Shadow PC first; if it persists, update the GPU
driver through Shadow's own updater or GeForce/NVIDIA Experience, then repeat this check.
\textbf{Do not continue until this works} --- without the driver, the benchmark will run
on CPU (uselessly slow) and the telemetry this trial exists to validate will be empty.

\section*{Part 2 --- Install Git for Windows (2 minutes)}

Git fetches the repository; the bundled \emph{Git Bash} shell is also what the
project's scripts (and Claude Code) run in.

\textbf{[PowerShell]}
\begin{cmdblock}
winget install --id Git.Git -e --source winget
\end{cmdblock}

Notes: \texttt{winget} is Windows' built-in package manager. If it reports
``\texttt{winget is not recognized}'', install ``App Installer'' from the Microsoft
Store, close and reopen PowerShell, and retry.

\textbf{Close PowerShell and open a new one} (installs only appear in fresh windows).

\checkbox \texttt{git --version} in the new PowerShell prints something like
\texttt{git version 2.5x}.

\section*{Part 3 --- Install Python 3.13 (2 minutes)}

\textbf{[PowerShell]}
\begin{cmdblock}
winget install --id Python.Python.3.13 -e --source winget
\end{cmdblock}

Again close PowerShell and open a fresh one.

\checkbox \texttt{python --version} prints \texttt{Python 3.13.x}.
(If it opens the Microsoft Store instead: Start $\rightarrow$ ``Manage app execution
aliases'' $\rightarrow$ turn \textbf{off} both \texttt{python.exe} and
\texttt{python3.exe} aliases, then retry in a fresh window.)

\section*{Part 4 --- Install Claude Code (3 minutes)}

Claude Code is the AI command-line agent that will drive the whole trial for you
(Part 6, Path A). If you prefer to run everything by hand, you may skip this part
and use Path B --- but Path A also handles surprises for you.

\textbf{[PowerShell]}
\begin{cmdblock}
irm https://claude.ai/install.ps1 | iex
\end{cmdblock}

Open a fresh PowerShell, run \cmd{claude}, and follow the login prompts in the
browser window it opens (log in with the Anthropic account). Exit Claude Code
afterwards by typing \texttt{/exit}.

\checkbox \texttt{claude --version} prints a version number without error.

\section*{Part 5 --- Get the repository (1 minute)}

The repository is public --- no login needed to download it.

\textbf{[PowerShell]}
\begin{cmdblock}
cd \textasciitilde \\
git clone https://github.com/ericmedlock/MSAI-LLM-PERFORMANCE.git
\end{cmdblock}

(\texttt{\textasciitilde} means your home folder, e.g.\ \texttt{C:\textbackslash Users\textbackslash you}.)

\checkbox A folder \texttt{MSAI-LLM-PERFORMANCE} now exists in your home folder and
contains \texttt{scripts}, \texttt{tasks}, \texttt{config}, \texttt{results},
\texttt{docs}.

\section*{Part 6 --- Run the trial}

\subsection*{Path A --- let Claude Code drive (recommended)}

Open \textbf{Git Bash} (Start key $\rightarrow$ type \texttt{git bash} $\rightarrow$ Enter), then:

\begin{cmdblock}
cd \textasciitilde/MSAI-LLM-PERFORMANCE \\
claude
\end{cmdblock}

When Claude Code starts, type exactly:

\begin{center}
\Large\texttt{RUN SHADOW PC TRIALS}
\end{center}

The repository ships a skill that tells Claude precisely what to do: run the setup
script, verify the GPU telemetry in the first rows, run the trial, then summarize
and commit the results. Approve the permission prompts it shows (they correspond to
the exact commands in Path B below). You can ask it questions at any time
(``how far along is it?'', ``show accuracy so far'').

\subsection*{Path B --- manual (identical result, two commands)}

Open \textbf{Git Bash}:

\begin{cmdblock}
cd \textasciitilde/MSAI-LLM-PERFORMANCE \\
bash scripts/setup.sh \\
bash scripts/run\_trials.sh
\end{cmdblock}

\subsection*{What the setup step does (either path), and what you'll see}

\begin{enumerate}
  \item Detects the platform and prints \texttt{[setup] detected profile: shadow}.
  \item Creates a private Python environment and installs the pinned dependencies
    ($\sim$2 min).
  \item Runs the offline test suite --- \checkbox expect
    \texttt{112 passed, 2 deselected}.
  \item Installs \textbf{Ollama} (the local model server) via winget --- expect a UAC
    prompt --- and starts it.
  \item Downloads the pinned model \texttt{deepseek-r1:14b} ($\sim$9 GB; Shadow's
    datacenter connection usually does this in a few minutes; a progress bar is shown).
  \item Prints the GPU line --- \checkbox expect \texttt{NVIDIA RTX A4500, 24564 MiB}
    and finally \texttt{[setup] DONE}.
\end{enumerate}

The run step then prints a dry-run plan --- \checkbox expect
\texttt{Total cells : 108 runs}, writing to
\texttt{results/frontier-v2-shadow-trial-14b.jsonl} --- and starts executing.
One line of progress appears per finished run.

\section*{Part 7 --- The five-minute telemetry check (the point of the trial)}

After the first 2--3 runs finish ($\sim$5 minutes in), open a \textbf{second} Git Bash
window and inspect the first recorded row:

\begin{cmdblock}
cd \textasciitilde/MSAI-LLM-PERFORMANCE \\
head -n 1 results/frontier-v2-shadow-trial-14b.jsonl \textbackslash \\
\mbox{\quad}\textbar{} ./.venv/Scripts/python -m json.tool \textbar{} grep -iE "vram\textbar gpu\textbar watt"
\end{cmdblock}
(The \texttt{\textbackslash} at the end of the first line lets you press Enter and
continue the command on the next line.)

\checkbox The GPU fields (\texttt{peak\_vram\_mb}, GPU utilization, power) contain
\textbf{non-zero numbers}. This proves the CUDA/NVML telemetry path works --- the
single most important output of this trial.

\textbf{If the fields are zero/empty:} stop the run (\texttt{Ctrl-C} in the first
window --- perfectly safe) and investigate before spending hours: is
\texttt{nvidia-smi} still working? Did Ollama load the model onto the GPU
(\cmd{ollama ps} should show the model with \texttt{100\% GPU})? Report what you find;
nothing collected so far is wasted.

\section*{Part 8 --- During the run}

\begin{itemize}
  \item \textbf{Duration:} roughly 2--3 hours for all 108 runs. Math items (AIME) are the
    slowest --- several minutes each is normal; some will legitimately fail (the tier is
    calibrated to $\sim$50\% single-pass difficulty). Failures are data, not problems.
  \item \textbf{Interrupting is always safe.} Every finished run is written to disk
    immediately. \texttt{Ctrl-C} stops it; re-running
    \cmd{bash scripts/run\_trials.sh} continues exactly where it left off.
  \item \textbf{Do not disconnect the Shadow client} (Part 0). If you must, reconnect
    later and simply re-run the command --- it resumes.
  \item Progress check at any time (second Git Bash window):
\begin{cmdblock}
wc -l results/frontier-v2-shadow-trial-14b.jsonl
\end{cmdblock}
    The number printed is finished runs out of 108.
\end{itemize}

\section*{Part 9 --- Send the data home}

Pushing requires GitHub authentication once. The easiest route is the GitHub CLI:

\textbf{[PowerShell]}
\begin{cmdblock}
winget install --id GitHub.cli -e --source winget
\end{cmdblock}

Then in a fresh \textbf{Git Bash}:
\begin{cmdblock}
gh auth login
\end{cmdblock}
(choose \texttt{GitHub.com} $\rightarrow$ \texttt{HTTPS} $\rightarrow$ login via
browser), and identify yourself to git (once):
\begin{cmdblock}
git config --global user.name "Eric Medlock" \\
git config --global user.email "ericmedlock@gmail.com"
\end{cmdblock}

Finally commit and push the results (Path A: Claude Code offers to do this for you):
\begin{cmdblock}
cd \textasciitilde/MSAI-LLM-PERFORMANCE \\
git add results/frontier-v2-shadow-trial-14b.jsonl \\
git add results/host/shadow.json results/hosts.csv \\
git commit -m "data(shadow): frontier-v2 trial N=1, RTX A4500" \\
git push
\end{cmdblock}

\checkbox \texttt{git push} ends with no error, and on the repository's GitHub
page the new file appears under \texttt{results/}.

\section*{Part 10 --- Troubleshooting quick reference}

\begin{tabular}{p{2.6in}p{3.4in}}
\textbf{Symptom} & \textbf{Fix}\\[2pt]
\hline\\[-8pt]
\texttt{winget} not recognized & Install ``App Installer'' from Microsoft Store; fresh PowerShell.\\[4pt]
\texttt{git}/\texttt{python}/\texttt{claude} not recognized right after install & Close the window; open a fresh one (PATH updates only apply to new windows).\\[4pt]
\texttt{python} opens the Microsoft Store & Start $\rightarrow$ ``Manage app execution aliases'' $\rightarrow$ disable the \texttt{python.exe} aliases.\\[4pt]
\texttt{ollama} not recognized after setup & Fresh Git Bash window; or run \texttt{bash scripts/setup.sh} again (it is safe to repeat).\\[4pt]
Model download very slow / stalled & \texttt{Ctrl-C} and re-run setup; Ollama resumes partial downloads.\\[4pt]
GPU fields empty in Part 7 & See Part 7; check \texttt{ollama ps} shows \texttt{100\% GPU}; check \texttt{nvidia-smi}.\\[4pt]
Run seems frozen on one item & AIME math items can take several minutes of silent ``thinking''; the harness also has its own timeout. Give it 10 minutes before worrying.\\[4pt]
Shadow disconnected overnight & Reconnect, open Git Bash, \texttt{cd} to the repo, re-run \texttt{bash scripts/run\_trials.sh} --- it resumes.\\[4pt]
\end{tabular}

\vspace{1em}
\section*{Final checklist}

\begin{itemize}
  \item \checkbox \texttt{nvidia-smi} shows the RTX A4500 \hfill (Part 1)
  \item \checkbox Git, Python 3.13, Claude Code installed \hfill (Parts 2--4)
  \item \checkbox Repo cloned \hfill (Part 5)
  \item \checkbox Setup finished: 112 tests passed, model pulled, \texttt{[setup] DONE} \hfill (Part 6)
  \item \checkbox First-row GPU telemetry is non-zero \hfill (Part 7)
  \item \checkbox 108/108 rows in the results file \hfill (Part 8)
  \item \checkbox Results committed and pushed to GitHub \hfill (Part 9)
\end{itemize}

\textit{Rules the scripts already enforce, stated for completeness: never edit
\texttt{config/config.yaml}, anything in \texttt{tasks/}, or \texttt{prompts/} ---
these are frozen by pre-registration, and the config file's bytes are hashed into
every result row. Machine-specific overrides belong in a \texttt{.env} file only.}

\end{document}
